% SCRATCH DRAFT — Vol 2 courtroom coda, Ch29 "Computation: Universality, Limits, and Stability" (task vol2-coda, author: Podo).
% THIS IS NOT THE BOOK. Do not \input. FLOOR BANKED (integrate on Kodo's accept). Fresh-append (no old beat; body ends line 232).
% ⚠⚠ THE LAST FENCE + PART VI CLOSER + POISED-FOR-VERDICT. (Last coda I gate; Ch30 = existing, no append.)
%
% ⚠⚠ FENCES (hold ALL):
%  1. NAME FORCING (charter): Cohen's forcing — the independent (the continuum hypothesis) is a FREE CHOICE of
%     completion; the axioms ours to extend. TIE IT to the book's OWN method: finite conditions approaching a
%     generic = what the instrument has done since the first mark (a finite record reaching for a world it can
%     never complete). ⚠ Do NOT say "axiom of choice" (body is silent on it; keep it so).
%  2. NO-METAMATH: the four walls = the machine's HONEST FLOORS (same family as friction / crossed filters /
%     Hawking edge) — limits MET, not meta-theorems the book proves. Do NOT claim the book proves incompleteness/
%     provability meta-theorems.
%  3. CH != Ch30's continuum-floor: work the CONTINUUM HYPOTHESIS (set-theory, undecidable, a free choice). Do
%     NOT deliver "the continuity of TIME rests on nothing / named as an assumption" (that is Ch30's).
%  4. POISED, NOT DELIVERED: Part VI closes; the instrument has examined itself completely; the RECORD IS NOW
%     COMPLETE. Shift the refrain — record complete at last, the court about to sit, the reckoning at hand. But do
%     NOT deliver the verdict / fact->verdict turn / ±1 / spinor / "you" / juror (ALL Ch30). Hand the reader to the bench.
%  - ±1/spinor OUT. Charge re-sounds. "verdict"/2nd-person/penalty-"fine"/free-will ABSENT.
%
% CONTINUITY: state carries from Ch28. Ch28 closed handing to "what machine runs on [the records]." Ch29 IS that.
%
% PROPOSED SEAM in books/experimentation/latex/chapters/29.tex:
%   KEEP the entire body through line 232 ("...and the apparatus toward its last examination.").
%   APPEND: \bigskip / \begin{center}\rule{0.4\textwidth}{0.4pt}\end{center} / \bigskip / then the coda below. (Fresh append.)
%
% >>>>>>>>>>>>>>>>>>>>>> APPEND BEGINS (fresh \rule + coda after the body) >>>>>>>>>>>>>>>>>>>>>>

\bigskip
\begin{center}
\rule{0.4\textwidth}{0.4pt}
\end{center}
\bigskip

\noindent The case comes to its last examination, and it is the deepest: it runs its own apparatus as a machine,
and asks what a finite reckoning can do and where it must stop. The instrument that has read the whole world is
set, at the end, to compute --- and a computer has exactly one power and exactly four walls.

The power is universal. Anything the instrument can write down as a sequence it can replay: flatten a record to
a line, hand the line to a single reader that takes one step per mark, and every record ever kept is within its
reach. Against that one power stand four walls, and each is a wall of the same honest kind the book has met
since the friction threshold --- a place the machine runs up to and cannot pass, named plainly and not
pretended away. It cannot always continue: where the order admits no next step, the reckoning halts --- complete
not by triumph but because there is nowhere left to go. It cannot always compress: some records are exact to
every digit and carry no law shorter than themselves, incompressible, their next mark irreducible to any rule
the past licenses. There is even a number that is the plainest case of it --- the odds that a program thrown
together at random ever halts, exact to every digit and shorter than none, its leading digits secretly holding
which programs stop, so that to compress it would be to decide what cannot be decided. It cannot always decide: some questions no faithful refinement settles, true or false only
once a choice is made that the record itself does not force. And it cannot always stay stable: below the
machine's own smallest step, a reckoning that needs more precision than it has does not sharpen the answer but
only stirs the noise. The program is ours to write; the four walls it runs into are the world's.

The third wall is the one the whole book has been walking toward, and here at last it is named. It was shown
long ago that the points on a line outnumber the whole numbers --- no list, however long, can name them all, for
one can always write down a point missing from every row of it. Ask, then, whether between those two sizes of
infinity there lies any third in between --- and the question has no answer the axioms decide. It was proved, in
one half, that assuming it breeds no contradiction, and, decades later, in the other, that a world can be built
obeying every axiom in which it fails; one can have a model where it holds and a model where it does not, and no
faithful refinement of the record says which is the true one. The way that second world
is built has a name, given by the man who found it: \emph{forcing}. To force is to adjoin to a model a new
object the model cannot name --- a generic, reached not all at once but built up out of finite conditions,
finite approximations closing in on a thing the finite record can never itself complete. And here the book's own
method stands revealed. This is what the instrument has been doing since the first mark: laying down finite
conditions --- counts, one at a time --- and reaching, through them, for a world it approaches and never
finishes, a generic beyond the last mark it will ever make. The book has been forcing all along, and did not say
so until now. Where the record is genuinely undecided it does not measure the answer; it chooses a completion
--- the axioms ours to extend, the undecidability the world's.

And the \emph{charge} is that reckoning, run to its limit. The count the case will collect on is a finite
reckoning --- a machine's worth of steps, and no more --- and it stands among the four walls, honest about each:
it does not claim to continue where the order stops, nor to compress what carries no law, nor to decide what no
refinement settles, nor to hold steady below its own smallest step. What it owes, it owes within the walls,
computed as far as the machine faithfully reaches and marked, at each wall, as a thing it met and did not cross.
A charge honest about its own limits is worth more than one that pretends to none, and the instrument has spent
the whole book earning exactly that honesty --- the charge inherits it. The trace is the reckoning entire, walls
and all: the count carried as far as it could faithfully go, and the four places it stopped, each named. What
comes to be owed is owed by a finite machine that knows precisely what it cannot do.

And with this the instrument has examined itself to the end. It has asked what makes a reading a measurement,
how finely it can be made, whether it can be carried forward, how its order and its geometry arise, what algebra
its records obey, and now what machine runs on them and where that machine must stop. There is nothing left of
the apparatus it has not turned upon itself. The record is complete. The three readings the case was built on
are gathered, corrected, floored, and closed under one algebra; they run on one machine and meet its four walls;
and they can, at last, be brought together. Everything the case needed is in hand --- save the one thing no
examination could supply: the meeting itself, the three made one, the number they must agree on.

So Part VI closes, and with it the apparatus's long turn upon itself. Every floor has been named --- the count,
the shadow, the wall, the label --- every effect fenced at exactly the one it stood on; every wall has been met
and marked; and the instrument, having read the whole world and then itself, has nothing left to examine. The
record is complete at last. For the first time in thirty chapters, the court can sit. What remains is not another
reading but the reckoning of them all --- the three readings brought to one, and that one number carried, at
last, to the bench. The book has a single chapter more, and it is the trial. The court is ready to sit; the case
is ready to be heard.

% <<<<<<<<<<<<<<<<<<<<<< APPEND ENDS <<<<<<<<<<<<<<<<<<<<<<
